\section{Related Work}
\label{sec:related}

\textbf{Protocol conformance and interoperability testing.}
Protocol conformance testing traditionally relies on curated test suites,
interoperability events, and implementation-specific regression tests. These
approaches are valuable because they produce executable evidence, but they
often cover only known behaviors and rarely maintain an explicit link between a
test and the normative sentence it checks. \tool{} complements conformance
testing by deriving tests from field-centric rules and by preserving the
standard-to-code trace for each report. Protocol state fuzzing learns TLS and
DTLS implementation automata and exposes deviations in deployed
stacks~\cite{de2015protocol,fiterau2020dtls}. Automata-based bug detection and
analyses of composite TLS state machines further show how protocol state
structure creates conformance risk~\cite{fiterau2023automata,beurdouche2015messy}.
Adversarial and combinatorial TLS testing frameworks further show how targeted
message generation can reveal library-specific behavior~\cite{somorovsky2016tlsattacker,maehren2022tlsanvil}.
Certificate-validation studies use generated, differential, and RFC-directed
inputs to find noncompliance in SSL/TLS implementations~\cite{brubaker2014frankencert,chen2015guided,chen2018rfc},
while hostname-verification and symbolic certificate analyses focus on
different slices of the same conformance problem~\cite{chen2017hvlearn,chau2017symcerts}.

An important difference is the direction of generation. Traditional conformance
suites are usually written by protocol experts who decide which cases are worth
testing. \tool{} starts from the standard and asks which implementation behavior
would falsify a field rule. This makes the generated tests less polished than a
hand-written suite, but it allows the system to surface gaps that existing tests
did not encode.

\textbf{Formal methods for protocol verification.}
Formal verification has produced strong results for cryptographic protocols,
state machines, and safety properties when a precise model is available. The
main obstacle for broad implementation checking is the cost of constructing and
maintaining such models from natural-language standards and production code.
\tool{} does not attempt to prove whole-protocol correctness. Instead, it
extracts lightweight field-level rules that can be aligned to code and checked
with a combination of static evidence and tests. This positioning is
complementary to verified TLS implementations such as
miTLS~\cite{bhargavan2013mitls}. It also complements general symbolic
verification frameworks, including ProVerif, Tamarin, and
AVISPA~\cite{blanchet2016proverif,meier2013tamarin,armando2005avispa}, as
well as symbolic analyses specialized to TLS 1.3~\cite{cremers2017tls13}.

This trade-off is deliberate. Full formalization provides stronger guarantees,
but it is rarely available for every implementation version and every extension.
Field-level rules are weaker but easier to derive, easier to inspect, and close
to the code locations where many implementation mistakes occur.

\textbf{Static analysis and fuzzing.}
Static analysis can find missing checks, unsafe data flows, and inconsistent
error handling through system-specific checkers and bug-pattern
inference~\cite{engler2001checking,engler2001bugs}. Software model checking
extends this style of reasoning to executable code paths~\cite{musuvathi2004cmc},
and security-oriented static analysis has scaled to both application code and
large deployed codebases~\cite{livshits2005finding,bessey2010few}. Fuzzing
exposes crashes and behavioral anomalies by generating structured or symbolic
inputs~\cite{godefroid2008grammar,godefroid2008sage,cho2011mace}, while
coverage-guided and targeted fuzzers improve exploration under finite
budgets~\cite{bohme2016aflfast,lemieux2018fairfuzz}. Protocol fuzzers adapt
these ideas to network implementations and stateful servers~\cite{pham2020aflnet,natella2022stateafl}.
Benchmarks and stateful greybox fuzzers further show the importance of
preserving protocol state while exploring message sequences~\cite{natella2021profuzzbench,ba2022stateful}.
Differential fuzzing provides another way to expose behavioral divergence
across implementations~\cite{petsios2017nezha}, and fuzzing surveys place
these techniques in a broader testing landscape~\cite{manes2021fuzzing}.
However, both techniques usually require either code-level specifications or
manually written oracles to decide whether an implementation violates a
standard. \tool{} uses field-centric rules as specification-derived oracles and
uses static analysis to ground LLM reasoning in real code structure.

The closest fuzzing systems infer message grammars or state machines and then
exercise implementations aggressively. \tool{} can use such generators as a
backend, but its novelty is the normative oracle: a generated input is valuable
because it distinguishes whether an implementation satisfies a specific rule,
not merely because it explores a new path.

\textbf{Specification mining and traceability.}
Prior work on specification mining infers likely rules from code, logs, tests,
or documentation. Dynamic invariant detection recovers behavioral properties
from executions~\cite{ernst2001dynamically}, while system-specific checkers and
deviant-behavior analyses infer likely rules from implementation
patterns~\cite{engler2001bugs,engler2001checking}. Requirements traceability
links natural-language requirements to source-code artifacts using information
retrieval and latent semantic indexing~\cite{antoniol2002recovering,marcus2003recovering}.
Candidate-link generation and broader traceability frameworks study how such
links should be produced, evaluated, and maintained~\cite{hayes2006advancing,clelandhuang2014software}.
More recent traceability work adds semantic representations for linking
requirements and code~\cite{guo2017semantically}, and protocol-specific
specification mining can extract finite-state models from standards
documents~\cite{pacheco2022automated}. \tool{} shares the goal of connecting
text and code but focuses on protocol standards, where fields provide a natural
semantic anchor and where conformance claims require both normative provenance
and behavioral evidence.

Unlike general traceability, a wrong link in this setting can become a false
vulnerability report. \tool{} therefore keeps provenance, confidence, and
uncertainty in the data model rather than treating trace recovery as a standalone
ranking problem.

\textbf{LLMs for software engineering.}
LLMs have been applied to code search, bug finding, patch generation, program
repair, and test generation. Directly applying LLMs to conformance checking is
appealing but unreliable because the model must resolve specification meaning,
identifier mapping, hidden context, and runtime behavior at once. \tool{}
separates these concerns: it constrains extraction with a field space, narrows
code evidence with syntax-aware retrieval, requires structured uncertainty, and
validates suspicious claims with tests. Our use of LLMs builds on code-focused
pretraining for joint program and natural-language
representations~\cite{feng2020codebert,guo2021graphcodebert}. It also draws
on program-synthesis work showing that large models can generate programs from
natural-language intent~\cite{chen2021codex,austin2021program}. Open
code-generation models further demonstrate this capability at larger scale and
across multi-turn settings~\cite{nijkamp2023codegen,li2023starcoder}.

This design also differs from LLM bug-finding tools that optimize for broad
coverage of likely defects. \tool{} optimizes for defensible conformance claims:
every report should be explainable in terms of a standard span, a field rule, a
code path, and an observed behavior.

\textbf{Agentic program analysis.}
Recent systems use agents to iteratively inspect repositories, run tools, and
generate tests. \tool{} is an instance of this broader direction, but its agent
loop is specialized around protocol conformance: missing evidence is expressed
as a request for implementation context relevant to a specific field rule, and
test generation is guided by counterexamples to a normative requirement.
ReAct-style agents interleave reasoning with external actions~\cite{yao2023react},
while Reflexion and Self-Refine illustrate feedback loops that revise an
agent's behavior or outputs across iterations~\cite{shinn2023reflexion,madaan2023selfrefine}.
